%! TEX root = main.tex

본 연구는 민감한 수치형 의료 데이터의 프라이버시를 완벽히 보호하면서도 실시간 검색을 가능하게 하는 영지식 증명(zk-SNARK) 기반의 새로운 사전 증명 검색 시스템을 제안하고 그 실효성을 입증하였다. 기존의 암호화 데이터베이스 연구들은 일부 정보 노출을 허용하여 보안성이 훼손되거나(OPE 등), 완전한 은닉을 추구하다가 실용성을 상실하는(ZKSQL 등) '데이터 프라이버시와 검색의 이중성(Dichotomy)'이라는 근본적인 딜레마에 빠져 있었다. 본 연구는 데이터의 원문을 서버에 일절 저장하지 않고 '사전 정의된 임상 참고치'에 대한 범위 증명 객체만을 저장하는 방식을 도입함으로써 이 문제를 성공적으로 타파하였다.

\subsection{연구의 의의}
본 연구의 가장 큰 학술적, 실용적 의의는 다음과 같다. 첫째, 기존 검증 가능한 데이터베이스 시스템의 치명적 한계였던 실시간 증명 연산 병목 현상을 완벽히 제거하였다. 설계 방식 비교 실험을 통해 제안 시스템이 응답 반환 개수(K)에 비례하여 선형적으로 지연되는 기존 모델의 한계(O(K))를 극복하고, 순수 데이터베이스 I/O 시간에 수렴하는 O(1) 스케일의 압도적인 검색 속도를 달성함을 확인하였다. 둘째, 다중 쿼리(Composite Query) 처리 시 필연적으로 발생하는 '소유권 증명' 문제를 단일 그룹 커밋과 실시간 증명의 일괄 생성을 통해 해결하였다. 이를 통해 다수의 서브 쿼리를 처리할 때도 암호학적 연산 지연을 500ms 미만으로 통제하였으며, 네트워크 부하 역시 약 4.9KB 수준으로 최소화하여 상용 클라우드 서비스 환경에서도 충분히 도입 가능한 수준의 실용성을 확보하였다.

\subsection{연구의 한계} 
그러나 본 연구는 대규모 검색 시스템으로서의 '실시간성'을 달성하기 위해 몇 가지 의도된 시스템적 절충을 채택하였으며, 이는 다음과 같은 한계점을 지닌다. 첫째, 다중 쿼리 결과의 실시간 그룹화를 위해 검색 서버(SS) 내부에 외부와 격리된 OWNERSHIP 매핑 테이블을 유지하였다. 비록 원문 데이터는 영지식 증명 뒤에 안전하게 은닉되지만, 서버가 완전히 장악당하는 최악의 시나리오에서는 특정 소유자의 데이터 분포 정보가 노출될 수 있는 위험이 존재한다. 둘째, 본 시스템은 서버를 '반신뢰(Semi-honest) 주체'로 가정하였다. 결과에 포함된 데이터들이 조작되지 않았음(Soundness)은 수학적으로 완벽히 증명하지만, 악의적인 서버가 조건에 부합하는 결과를 고의로 누락하는 행위(Incompleteness)를 원천적으로 검증하거나 방어하지는 못한다. 이는 완전한 은닉성 하에서 무결성을 검증하기 위해 막대한 통신 및 연산 오버헤드를 유발하는 기존 연구들의 한계를 피하기 위한 불가피한 선택이었다.

\subsection{향후 과제}
본 연구의 한계를 보완하고 시스템의 범용성을 넓히기 위해 다음과 같은 후속 연구를 제안한다. 첫째, 악의적 위협 모델 하에서도 검색의 실시간성을 크게 훼손하지 않으며 쿼리 결과의 누락 없음(Completeness)을 증명할 수 있는 고도화된 검증 프로토콜에 대한 연구가 필요하다. 둘째, 현재 시스템은 사전 정의된 이산적인 구간(Pre-defined Range) 정보에 기반하고 있다. 향후에는 동형 암호나 다자간 컴퓨팅(MPC) 등의 최신 암호학적 원형들을 영지식 증명과 결합하여, 연속적인 범위 검색이나 텍스트와 같은 비수치형 데이터에 대해서도 온디맨드(On-demand)로 실시간 증명이 가능한 범용적인 프라이버시 보존 검색 아키텍처로 발전시켜 나갈 계획이다.